\section{The Foundation Triad}
\label{sec:foundation}

\subsection{KFD-1: continuity of the fact world}

KFD-1 states that facts must not drift. A \emph{fact} here is not whatever a
system happens to store, nor a claim of observer-free certainty. It is a
load-bearing reference admitted from a declared source, authority, and
evidence boundary. KFD distinguishes a \emph{fact cut}, which states what
holds at such a boundary, from a \emph{causal record}, which states what
occurred between declared boundaries \cite{kfd2026alpha42}. Either may include
known gaps, uncertainty, or conflicting evidence.

Let $F_t$ be a fact cut used at time $t$, and let $c$ be its declared
compatibility boundary. If the system presents the same fact identity and
boundary at $t'$ without declaring a change, its load-bearing meaning must be
preserved:

\begin{equation}
\operatorname{id}(F_t)=\operatorname{id}(F_{t'}) \land c_t=c_{t'}
\Longrightarrow
\operatorname{meaning}(F_t)=\operatorname{meaning}(F_{t'}).
\label{eq:fact-continuity}
\end{equation}

The invariant does not freeze the world. It requires semantic change to cross
an explicit boundary so later evidence can contradict, supersede, or narrow an
earlier model without silently rewriting what the earlier participant saw.

KFD-1 therefore provides \emph{world continuity}. It lets two participants
refer to the same admitted state or occurrence even when they disagree about
its interpretation. Without such a reference, disagreement cannot reliably
be distinguished from drift in the object being discussed.

\subsection{KFD-2: warranted reliance}

KFD-2 states that trust must start from facts. Trust is not modeled as a
global property of a person, model, product, or artifact. It is a bounded
permission to rely on a claim for a purpose.

For claim $C$ and purpose $P$, a trust assessment can be represented as:

\begin{equation}
T_P(C)=\mathcal{A}(C,P,F,E,R,Q),
\label{eq:trust-assessment}
\end{equation}

where $F$ is the bound fact cut, $E$ the checked evidence, $R$ residual risk,
and $Q$ the assignment of source, verification, and decision responsibility.
The result may be pass, warning, fail, or unverifiable. Its force is limited to
the stated purpose and cut.

This definition explains why more data is not the same as more trust. A log
may be authentic while irrelevant to a completion claim. A check may pass
while omitting the platform a consumer uses. A capable agent may produce a
persuasive explanation without exposing the facts that would let another
participant challenge it. KFD-2 makes reliance revisable because contrary
facts can change the assessment without first defeating the authority or
reputation of the claimant.

KFD-2 therefore provides \emph{judgment continuity}. It allows a later
participant to recover not only what was claimed, but why and within what
boundary the claim was permitted to guide action.

\subsection{KFD-3: generative coordination}

KFD-3 states that cooperation must start from trusted value. Facts and trust
do not themselves create joint action. A system may expose accurate evidence
while hiding the value exchange, available alternatives, or constraints under
which participation is requested.

Let $V$ be a participant-relevant value claim, $T(V)$ its trust assessment,
$K$ the visible constraints, and $O$ the meaningful available options. A
cooperation decision has the form:

\begin{equation}
C = \mathcal{C}(V,T(V),K,O,Q_C),
\label{eq:cooperation}
\end{equation}

where $Q_C$ declares responsibility for the commitment and resulting action.
The participant need not agree with every other participant. It must be able
to understand what is offered, what is constrained, and what accepting,
rejecting, escalating, or withdrawing means within the declared boundary.

This is not an absence-of-control rule. Sandboxes, permissions, refusal,
revocation, and hard safety gates may be necessary. KFD-3 distinguishes
explainable, reviewable constraint from hidden capture, and cooperation from
mere compliance. The distinction matters operationally: a compliant component
executes inside one authority; a cooperating participant can contribute new
facts, objections, alternatives, and commitments to a shared world.

KFD-3 therefore provides \emph{action continuity}. It lets independent
participants extend the work without surrendering the fact and judgment
boundaries that made the extension trustworthy.

\subsection{Three irreducible questions}

The triad can be summarized by three questions:

\begin{table}[ht]
\centering
\small
\begin{tabular}{p{0.15\linewidth}p{0.24\linewidth}p{0.48\linewidth}}
\toprule
Layer & Role & Question preserved across participants and time \\
\midrule
KFD-1 & World continuity & What state or occurrence are we talking about, and
where may its meaning change? \\
KFD-2 & Judgment continuity & Why may this claim be relied upon for this
purpose, and what remains unproved? \\
KFD-3 & Action continuity & Why and within what visible choices and constraints
should independent participants act together? \\
\bottomrule
\end{tabular}
\caption{The three operational roles of the Foundation Triad.}
\label{tab:triad-roles}
\end{table}

The roles correspond to reality as admitted, belief as warranted, and action
as coordinated. They are related but not interchangeable. That separation is
what makes the foundation testable.
